\section{電子計算機のイロハ}
\label{interMission}

\subsection{授業内容}
\subsubsection{科目の中でのこのコマの位置づけ}
\begin{tabular}[t]{ccccccc}
    記述統計[2] & 確率[0] & モデル[0] & 推測・推定[0] & 意思決定[0] & 実践[0]
\end{tabular}
\subsubsection{概要}
昨今はスマートフォン，タブレットなどのデバイスが行き渡っているが，専門的な計算をするためにはやはりパソコンを使うことになる。この講義は次回からのPCを使った演習に対する準備的内容であり，必ずしも心理統計に直結する内容ではない。とはいえ，統計計算の歴史は計算機発展の歴史でもあり，実際にPCのもつ基本的な知識・概念・技能がないと以降のPCを使った演習にも差し障りが出るので，十分な準備をしてから実践に進む必要があるだろう。
\subsubsection{コマ主題細目(箇条書き)}
\begin{description}
    \item[コンピュータの基礎]
    \item[情報の単位]
    \item[ファイル・フォルダ・プログラム]
    \item[再現性のある研究実践へ]
\end{description}
\subsubsection{キーワード}
\begin{itemize}
\item コンピュータの5大機能
\item ビット，バイト
\item 文字コード
\item フォルダと拡張子
\item アルゴリズム
\end{itemize}
\subsection{授業情報}
\paragraph{コマの展開方法}
講義
\paragraph{標準シラバスにおける位置づけ}
科目番号5；心理学統計法；2.統計に関する基礎的な知識；C エクセル，R，SPSS入門(１) 統計分析のための言語の手ほどき　プログラムの基本的操作
\subsubsection{予習・復習課題}
\paragraph{予習}
\paragraph{復習}
